\titleformat{\chapter}[display]
  {\normalfont\bfseries}{}{0pt}{\Large} % Este comando es para eliminar la aparición de "Capítulo N" en el PDF. Para emplear el formato de la tesis simplemente eliminamos la línea. 
\chapter{Materiales y métodos}
\section{Construcción y optimización del sistema C\textsubscript{60}(Gly)\textsubscript{4}}
La construcción de la estructura tridimensional del aducto C\textsubscript{60}(Gly)\textsubscript{4} se llevó a cabo utilizando el visualizador y editor molecular GaussView, tomando como referencia la metodología reportada por Sharoyko y colaboradores. La optimización de las geometrías moleculares se realizó mediante la teoría del funcional de la densidad (DFT) empleando el nivel de teoría B3LYP/6-311G(d,p). Para confirmar que las estructuras optimizadas corresponden a un mínimo local real en la superficie de energía potencial (SEP), se realizaron los cálculos de frecuencias vibracionales, verificando la ausencia de frecuencias imaginarias.
\section{Determinación de la reactividad local}
Para determinar los sitios preferenciales de adición para un ataque nucleofílico del grupo amino de la glicina sobre la jaula de C\textsubscript{60}, se calcularon las funciones de Fukui (FF) como descriptores de reactividad local. Dado que la glicina actúa como especie nucleofílica a través de su grupo amino y la jaula de C\textsubscript{60} se comporta como el electrófilo, se evaluó la función de Fukui para ataques nucleofílicos ($f^+$). 

La reactividad local se describe mediante las funciones de Fukui. En su representación espacial, la aproximación por diferencias finitas se expresa de la siguiente manera:

\begin{equation}
    \text{Ataque nucleofílico: } f^+(\mathbf{r}) = \rho_{N+1}(\mathbf{r}) - \rho_N(\mathbf{r})
\end{equation}

\begin{equation}
    \text{Ataque electrofílico: } f^-(\mathbf{r}) = \rho_N(\mathbf{r}) - \rho_{N-1}(\mathbf{r})
\end{equation}

\begin{equation}
    \text{Ataque radicalario: } f^0(\mathbf{r}) = \frac{1}{2} \left[ \rho_{N+1}(\mathbf{r}) - \rho_{N-1}(\mathbf{r}) \right]
\end{equation}

En estas expresiones, $\rho_{N+1}(\mathbf{r})$, $\rho_N(\mathbf{r})$ y $\rho_{N-1}(\mathbf{r})$ representan la densidad electrónica en la posición espacial $\mathbf{r}$ para las especies en sus estados aniónico ($N+1$), neutro (N) y catiónico ($N-1$), respectivamente. 

El procedimiento de adición del aminoácido sobre la jaula de C\textsubscript{60} se llevó a cabo de manera iterativa. Se identificaron los átomos de carbono con los mayores valores del índice $f^+$, y se adicionó la siguiente molécula de glicina en el carbono con la reactividad más pronunciada para posteriormente optimizar la geometría, repitiendo este protocolo para cada etapa de adición de un nuevo aminoácido.

Las propuestas estructurales se optimizaron mediante la teoría del funcional de la densidad (DFT) empleando el nivel de teoría B3LYP/6-311G(d,p), y se corroboró la ausencia de frecuencias imaginarias mediante el análisis de frecuencias vibracionales para asegurar que correspondieran a mínimos locales en la superficie de energía potencial. La comparación de las propuestas estructurales se realizó mediante sus diferencias de energía relativa. Las configuraciones evaluadas y las seleccionadas para continuar la funcionalización se presentan en Resultados, hasta obtener el aducto C\textsubscript{60}(Gly)\textsubscript{4}.
\section{Evaluación de la reactividad global y mapa de potencial electrostático}
Se evaluaron los descriptores de reactividad global para analizar la facilidad del aducto para donar o aceptar electrones y su relación con un posible comportamiento antioxidante. Para ello, se calcularon las energías del sistema en sus estados neutro, aniónico y catiónico, lo que permitió obtener el potencial de ionización ($I$) y la afinidad electrónica ($A$). Estos parámetros permitieron caracterizar la susceptibilidad y predisposición del sistema a reaccionar ante diversas especies químicas.

Para identificar los sitios propensos a interacciones electrostáticas, se generó la superficie de mapa de potencial electrostático (MEP, por sus siglas en inglés) para el ligando C\textsubscript{60}(Gly)\textsubscript{4}. El MEP evalúa la distribución de carga eléctrica sobre una superficie de densidad electrónica constante alrededor de la molécula. La visualización se realizó mediante un código de colores estándar: las regiones rojas representan potencial negativo (zonas ricas en electrones o de carácter nucleofílico), las regiones azules representan potencial positivo (zonas deficientes en electrones o de carácter electrófilo) y las regiones verde-amarillas indican zonas con potencial cercano a cero en la escala de representación utilizada. De esta manera, se determinaron los sitios clave de la molécula capaces de establecer interacciones por fuerzas electrostáticas con el entorno proteico.

\section{Acoplamiento molecular}
Para explorar el potencial inhibidor del sistema de fulereno C\textsubscript{60} funcionalizado con 4 glicinas sobre la proteína de replicación del virus del dengue (complejo NS2B-NS3) en su conformación abierta e inactiva, se llevó a cabo la reconstrucción y completación de la estructura utilizando la plataforma CHARMM-GUI. A partir de la estructura cristalográfica obtenida del Protein Data Bank (PDB ID: 2FOM), se modelaron los residuos no definidos correspondientes a los extremos N- y C-terminales. La calidad geométrica del modelo se evaluó mediante un mapa de Ramachandran, verificando los ángulos de torsión permitidos ($\phi$ y $\psi$) en la cadena principal. Posteriormente, la estructura se sometió a una minimización de energía empleando el campo de fuerza GROMOS 54A7 en el programa GROMACS.

El modelo resultante se preparó para el acoplamiento molecular a través de la interfaz AutoDockTools. Durante este proceso, se eliminaron las moléculas de agua cristalográficas, se añadieron los hidrógenos polares necesarios y se asignaron las cargas parciales de tipo Kollman. Asimismo, se revisaron manualmente los residuos sin carga asignada para corregir posibles desbalances energéticos antes de exportar la estructura en formato PDBQT.

Finalmente, se definió una caja de grilla (\textit{grid box}) centrada en la región de interés con los parámetros ajustados para la simulación. Para el análisis posterior de las conformaciones, se estableció un umbral de tolerancia de RMSD de 8 Å para el agrupamiento de las poses, determinado a partir de la visualización estructural del complejo proteína-ligando.
\section{Simulaciones de acoplamiento molecular con AutoDock Vina}
Las simulaciones de acoplamiento molecular (\textit{docking}) se llevaron a cabo utilizando el software AutoDock Vina. Las estructuras del receptor y los ligandos se prepararon y definieron en formato PDBQT. La ejecución del programa generó un total de 100 poses conformacionales, cada una asociada a su respectiva energía de afinidad de unión. Para la separación y gestión individual de estas conformaciones, se utilizó la herramienta \texttt{vina\_split}.

A partir de las conformaciones obtenidas, se realizó un análisis visual y un agrupamiento (\textit{clustering}) en función de su energía relativa, aplicando un umbral de corte basado en la raíz de la raíz de la desviación cuadrática media (RMSD). Finalmente, se construyeron los histogramas de distribución para evaluar las poses comprendidas dentro de dicho umbral de tolerancia en concordancia con sus correspondientes valores de energía de unión.


%La estructura tridimensional de la proteina (2FOM) se obtuvo por medio de la base de datos protein data bank. La generacion de los parametros de topología se llevarona cabo mediante el software Gromacs(), empleando el campo de fuerza GROMOS96 53A6.
% BORRADOR PREVIO: conservado para cotejo; contiene parámetros contradictorios
% con los apartados siguientes y campos sin completar.
% Construccion del entorno de misulacion y minimización de energía 
% La proteina se ubico en el centro de una caja de simulacion respetando una distancia minina de () nm entre la superficie de la proteina y el limite de la caja. El entorno se solvató utilizando el modelo de agua () para posteriormente añadir contraiones neutralizantes Cl- o Na+. A continuacuib, se ejecuto una minimización de energía mediante el algoritmo de (steepest descent)
% %Fase de equilibrio 
% El sistema se equilibro en dos estapas bajo ciertas restricciones: Primero una fase (nvt) empleando el termostato V-rescale/Berendsen donde el volumen es constante a 300K durante 10 ns, este  y el segundo una fase en el conjunto isotérmico-isobárico  (npt) donde la presion es constante a 300K durante 10 ns utilizando el barostato Parrinello-Rahman., el cual ajusta dinamicamente el tamaño y la forma de la caja de simulación en respeusta a las fluctuaciones de la presion. Este mecanismo permite que el volumen del sistema se expanda o contraiga de manera anisotropica hasta alcanzar una presion promedio (1atm), simulando asi un entorno mas realista y flexible, como el que se encuetra en los sistemas biologicos. 
% 
% Una vez completadas las fases de equilibrio, se procedio con la simulación de producción por un total de 100ns, empleando un paso de integracion de 0.0015ps y un total de 66,666,6667 pasos. Durante esta etapa, los enlaces covalentes que involucran atomos de hidrogeno se restringieron mediante el algoritmo LIBCS, el cual permite mantener fijas estas distancias sin afectar la dinámica general del sistema. 
% 
% Se almacenaron las trayectorias y coordenadas para su posterior análisis estructural. Asimismo, se eliminaron las moleculas de agua y los iones para simplificar el análisis de la trayectoria. Posteriormente, se calcularon los parametros estrcuturales de RMSD para evaluar la estabilidad dconformacional global, el RMSF para identificar regiones con mayor flexibilidad y el area superficial accesible al solvente  (SASA, Solvent Accessible Surfaces Area). aunado a ello se determino el radio de giro (Rg) como medida de compactación estructural y el numero de enlaces de hidógeno. Los resultados obtenidos se graficaron en (). 
% 
% %

\section{Preparación del sistema y solvatación}

La estructura tridimensional de la proteína (PDB ID: 2FOM) se obtuvo de la base de datos Protein Data Bank. La generación de los parámetros de topología y la preparación general del sistema se llevaron a cabo mediante el programa GROMACS. La proteína se ubicó en el centro de una caja de simulación cúbica, manteniendo una distancia mínima de 1.0 nm entre la superficie proteica y los bordes de la caja. Posteriormente, el entorno se solvató utilizando el modelo de agua SPC/E y se neutralizó la carga total del sistema mediante la adición de contraiones de Na\textsuperscript{+} y Cl\textsuperscript{-} empleando la herramienta gmx genion.

\section{Minimización de energía}

Para eliminar posibles choques estéricos y tensiones estructurales entre la proteína y el disolvente, el sistema se sometió a una minimización de energía utilizando el algoritmo de descenso más pronunciado (steepest descent). Esta etapa se configuró con un tamaño de paso inicial de 0.01 nm y un límite máximo de 50,000 pasos, dándose por concluida al alcanzar el criterio de convergencia con una fuerza máxima inferior a 1000.0 kJ/(mol·nm).

\section{Fases de equilibrio (NVT y NPT)}

El equilibrio del sistema se efectuó en dos etapas sucesivas manteniendo restricciones de posición sobre los átomos pesados de la proteína (-DPOSRES). En primer lugar, se llevó a cabo una fase NVT a volumen constante durante 1000 ps (1 ns) empleando el integrador leap-frog con un paso de tiempo de 2 fs ($dt = 0.002\text{ ps}$). La temperatura se reguló a 300 K mediante el termostato V-rescale ($\tau_t = 0.1\text{ ps}$), acoplando de manera independiente la proteína y el entorno de solvatación, mientras que las velocidades atómicas iniciales se asignaron mediante una distribución de Maxwell-Boltzmann.
A continuación, se realizó el equilibrio NPT a presión constante durante 1000 ps (1 ns), continuando directamente desde la fase previa. La presión se reguló de manera isotrópica a 1.0 bar mediante el barostato de Parrinello-Rahman ($\tau_p = 2.0\text{ ps}$) con un valor de compresibilidad de $4.5 \times 10^{-5}\text{ bar}^{-1}$, correspondiente al agua líquida. En ambas fases de equilibrio, las interacciones electrostáticas de largo alcance se trataron con el método Particle Mesh Ewald (PME) con un radio de corte de 1.0 nm y un espaciado de malla de Fourier de 0.16 nm, mientras que las interacciones de van der Waals se evaluaron con un límite de corte de 1.0 nm.

\section{Simulación de producción}

Una vez completado el equilibrio, se retiraron las restricciones de posición para proceder con la simulación de producción de dinámica molecular no restringida durante 100 ns ($50 \times 10^6$ pasos) bajo el ensamble NPT, utilizando un paso de integración de 2 fs ($dt = 0.002\text{ ps}$). La temperatura se mantuvo a 300 K mediante el termostato de Nosé-Hoover ($\tau_t = 1.0\text{ ps}$) y la presión a 1.0 bar con el barostato de Parrinello-Rahman ($\tau_p = 5.0\text{ ps}$). Para esta etapa, los cortes electrostáticos y de van der Waals se establecieron en 1.2 nm, aplicando una función de suavizado por fuerza (force-switch) a partir de 1.0 nm para van der Waals y un espaciado de Fourier de 0.12 nm para PME. Se incluyeron correcciones de dispersión para energía y presión, los movimientos de traslación del centro de masa se removieron periódicamente y todos los enlaces que involucran átomos de hidrógeno se restringieron mediante el algoritmo LINCS. Las coordenadas comprimidas de la trayectoria y los datos de energía se almacenaron cada 2 ps para su posterior análisis.

\section{Procesamiento de trayectoria y análisis estructural}

Previo al análisis estructural, las trayectorias se procesaron utilizando gmx trjconv para corregir los efectos de las condiciones periódicas de frontera (-pbc mol -center) y eliminar los movimientos globales de rotación y traslación mediante una superposición de mínimos cuadrados (-fit rot+trans). Asimismo, las moléculas de agua y los iones se removieron del sistema para simplificar el análisis. A partir de la trayectoria resultante, se evaluaron la raíz de la desviación cuadrática media (RMSD), la raíz de la fluctuación cuadrática media (RMSF), el radio de giro, el área superficial accesible al disolvente (SASA) y el número de puentes de hidrógeno. 








